%% kapa: exam=1 round=35 year=2024 period=2 form=A subject=appraisal_law subject_ko=감정평가관계법규 seq=7 number=7
%% source: Q-Net 감정평가사 기출문제 (cst00309) artlSeq=5213229 file=raw/1008/5213229/2교시 문제지 원본.pdf
%% conversion: kordoc 4.24.51 -> tools/gen1.py (auto); review=proofread-sonnet; flags=none
\begin{question}{7}
국토의 계획 및 이용에 관한 법령상 용도지역별 건폐율의 최대한도가 큰 순서대로 나열된 것은? (단, 조례 및 기타 강화ㆍ완화조건은 고려하지 않음)
\begin{framedtext}
\begin{tabular}{@{}p{0.45\linewidth}p{0.45\linewidth}@{}}
ㄱ. 제2종전용주거지역 & ㄴ. 유통상업지역 \tabularnewline
ㄷ. 일반공업지역 & ㄹ. 농림지역 \tabularnewline
\end{tabular}
\end{framedtext}
\begin{choices}
  \item ㄴ - ㄱ - ㄷ - ㄹ
  \item ㄴ - ㄷ - ㄱ - ㄹ
  \item ㄷ - ㄴ - ㄹ - ㄱ
  \item ㄷ - ㄹ - ㄱ - ㄴ
  \item ㄹ - ㄷ - ㄴ - ㄱ
\end{choices}
\end{question}
\answer{2}
